% Transcription of folder 147, batch 05 (archivists' pages 81-100).
% Pass: Opus 5.5 (claude-opus-5-5), 2026-10-09 — first pass, unchecked against the pages by a human.
%
% Grothendieck numbers every other sheet (author's pp. 40-49 on archivists'
% pp. 81, 83, ..., 99); the sheets in between continue the text. His
% equation numbers run from (105) to (133). Pages 82 and the upper half of 83
% are cancelled by diagonal strokes and are transcribed as such.
\documentclass[11pt,a4paper]{article}
\input{../preamble/grothendieck}

\folder{147}
\batch{5}
\pages{81}{100}
\foldertitle{Structure à l'infini des Mg,ν [pages 1 à 90, dont table des matières] : notes manuscrites (s.d.).}
\dating{s.d. — le groupe « Autour de La "Longue Marche" à travers la théorie de Galois » (141 à 149) est daté [à partir de 1978]-1983}
\watermark{Édition de démonstration}

\begin{document}

\page{81}\note{p. 40 de l'auteur (numéro cerclé). Le diagramme (105) qui ouvre la page fait suite à un argument commencé avant ce lot ; les renvois (103) visent une page antérieure.}
\note{Diagramme (105), redessiné. Dans plusieurs indices, un troisième indice a été biffé et est illisible ($\Pi M^{*}_{G,\cdot,G''}$ en haut à gauche, $\Pi M_{G,\cdot}$ en bas à gauche) ; la flèche $\gamma_{G,G'}$ aboutissait d'abord à un terme biffé illisible, suivi d'une flèche « can » vers $\Pi M_{G'}$. Le carré central porte la mention « (2 cart) ». L'étiquette de la flèche courbe de gauche est lue $\varphi_{G,G''}$ avec doute. Deux accolades renvoient à (103), celle de droite « (103) pour $(G',G'')$ ».}

\begin{tikzcd}[column sep=small, row sep=normal, nodes={font=\scriptsize}]
\Pi M^{*}_{G,G''} \arrow[rr, dashed] \arrow[d, dashed] \arrow[dd, bend right=40, "{\varphi_{G,G''}}"'] \arrow[rrr, bend left=20, "{\gamma_{G,G''}}"] & & \Pi M^{*}_{G',G''} \arrow[r, "{\gamma_{G',G''}}"] \arrow[d] & \Pi M_{G''} \\
\Pi M^{*}_{\underline{G},\underline{G}'} \arrow[rr, "{\gamma_{G,G'}}"] \arrow[d] & & \Pi M_{G'} & \\
\Pi M_{G} & & &
\end{tikzcd}

\[
\Pi M_{G'} \;\bigl(\simeq \Pi(M_{G'},\Gamma^{!}_{G',G''})\bigr)
\]

Dans ce diagramme, on a \struck{un peu} laissé tomber les indices pour les MD-graphes, \struck{et on a remplacé la notation}

\struck{$\Pi(M_{G,G'},\Gamma^{!}_{G,G'})$ par la notation plus légère $\Pi M_{G,G'}$, de même}

\struck{pour $\Pi M_{G',G''}$.} On \struck{lira} \struck{a défini} \uncertain{retrouve}
\[
(106)\qquad \Pi M^{*}_{G,G',G''}
\]
\note{le dernier indice de (106) est surchargé ; on lit $G''$ sous la rature.}
dans (105) comme un produit 2-cartésien -- c'est donc \ill{} \uncertain{groupoïde} de noyau\supplied{x} \ill{} cat.\ de $\Pi M^{*}_{G,G'}$ par un $\mathbb{Z}^{C'}$ (tordu), et il se trouve que c'est aussi, \uncertain{en tant que} gr.-cat.\ de $\Pi M_{G}$ (de noyau une cat.\ de $\mathbb{Z}^{C}$ tordu par $\mathbb{Z}^{C'}$ tordu) une extension de $\Pi M_{G}$ par un $\mathbb{Z}^{C\amalg C'}$ tordu. De façon plus précise, on retire de (105) le diagramme
\note{dans (106), (107) et (108), l'indice qui suit $G$ dans $\Pi M_{G,\cdot}$ est biffé et illisible.}

\begin{tikzcd}
\Pi M^{*}_{G,G',G''} \arrow[r, "{\gamma_{G,G',G''}}"] \arrow[d, "{\varphi_{G,G''}}"] & \Pi M_{G''} \\
\Pi M_{G} &
\end{tikzcd}

\note{diagramme numéroté (107).}

\page{82}\note{toute la page est barrée de deux grands traits en croix ; elle est transcrite telle qu'on la lit.}
\struck{que je voudrais comparer au diagramme de} \struck{généralisation} \struck{pour $G \to G''$ :}

\begin{tikzcd}
\Pi M^{*}_{G,G''} \arrow[r, "{\gamma_{G'',G}}"] \arrow[d, "{\varphi_{G,G''}}"] & \Pi M_{G''} \\
\Pi M_{G} &
\end{tikzcd}

\note{diagramme numéroté (108) ; l'étiquette de la flèche horizontale est bien $\gamma_{G'',G}$ (sic), et l'indice de $\Pi M_{G,\cdot}$ en bas est surchargé.}

\struck{Pour ceci, notons qu'on a un morphisme canonique}
\[
(109)\qquad \text{\struck{$\Pi M_{G,G'} \to \Pi M_{G,G''}$,}}
\]
\struck{qui provient du} \struck{morphisme} \add{\struck{dirigé, soit $\mathrm{W}$-morphisme}} \struck{canonique de germes de multiplicités $M_{G,G'} \to M_{G,G''}$), soit, un \ill{}, un \uncertain{relèvement} \ill{} --}
\[
(110)\qquad \text{\struck{$\Gamma^{!}_{G,G'} \subset \Gamma^{!}_{G,G''}$}}
\]
\struck{d'où}
\[
(111)\qquad \text{\struck{$\Pi(M_{G},\Gamma^{!}_{G,G'}) \xrightarrow{\ \mathrm{can}\ } \Pi(M_{G},\Gamma^{!}_{G,G''})$}}
\]
\note{sous chacun des deux termes de (111), un signe $\wr$ renvoie respectivement à $M_{G,G'}$ et à $M_{G,G''}$ ; le numéro (110) est lu sur un « (140) » douteux.}

\struck{Ceci posé, \struck{je dis} on a un diagramme \struck{qui} est commutatif}

\page{83}\note{p. 41 de l'auteur (numéro cerclé). La moitié supérieure de la page, jusqu'à la formule $\Pi M_{G,G'} \to \Pi M_{G,G''}$, est barrée de deux traits obliques.}

\begin{tikzcd}[column sep=small]
\Pi M_{G,G',G''} \arrow[rrr, "{\gamma_{G,G',G''}}"] \arrow[dr] \arrow[dd] & & & \Pi M_{G''} \arrow[dr] & \\
 & \Pi M^{*}_{G,G''} \arrow[rrr, "{\gamma_{G,G''}}"] \arrow[dd] & & & \Pi M_{G''} \\
\Pi M_{G} \arrow[dr, "\mathrm{can}"'] & & & & \\
 & \Pi M_{G,G''} & & &
\end{tikzcd}

\note{diagramme (112), barré ; l'indice de $\Pi M_{G,\cdot}$ à gauche est surchargé, et l'exposant $*$ du coin supérieur gauche est incertain.}

\struck{où le carré vertical est 2-cartésien. En d'autres termes, le produit 2-cartésien $\Pi M^{*}_{G,G',G''}$ défini formellement via (105)} \struck{comme composé de deux extensions} \struck{successives de $\Pi M_{G,G'}$, s'interprète aussi} \struck{directement comme l'extension \uncertain{image} inverse de $\Pi M^{*}_{G,G''}$ sur $\Pi M_{G,G''}$ par}
\[
\text{\struck{$\Pi M_{G,G'} \to \Pi M_{G,G''}$.}}
\]

Considérons maintenant le cas d'un morphisme de généralisation $G \to G'$, \supplied{où} $G'$ est le graphe de MD \emph{discret} défini par $G$ -- c'est donc la \add{généralisation} \struck{spécial\ill{}} « la plus grossière » de toutes. On \add{pose} \struck{désigne}

\struck{par $S\Pi M_{G}$}

\page{84}
\[
(113)\qquad S_{A}\Pi M_{G} \overset{\mathrm{déf}}{=} \Pi M^{*}_{G,G'} ,
\]
\marginal{ici $A$ désigne l'ens.\ des \struck{\ill{}} arêtes de $G$.}
\marginal{\struck{\ill{}} $= \surd$}
c'est donc une extension de $\Pi M_{G,G'}$ par $\mathbb{Z}^{A(G)}$ \struck{\ill{}}, où $A(G)$ est l'ensemble des arêtes de $G$.

Bien entendu, le groupe $\Gamma = \mathrm{Aut}\, G$ opère non seulement sur $M_{G}$ et sur $\Pi M_{G}$, mais aussi sur $S_{A}\Pi M_{G}$ \struck{\ill{}} (via son action sur $M^{*}_{G,G'}$), de façon à commuter \supplied{avec} les morphismes canoniques
\[
(114)\qquad S_{A}\Pi M_{G} \longrightarrow (\Pi M_{G},\ \cdot\,)
\]
\note{le second terme du couple, à droite de (114), est biffé et illisible.}
ce qui permet, pour \struck{$A$} \add{un} sous-groupe
\[
\Gamma' \subset \Gamma_{G}
\]
\struck{\ill{}} de définir une extension
\[
(115)\qquad (S_{A}\Pi M_{G},\Gamma') \longrightarrow (\Pi M_{G},\Gamma')
\]
\marginal{explicitons un peu…}
\[
\simeq S_{A}\Pi(M_{G},\Gamma') \longrightarrow \Pi(M_{G},\Gamma')
\]
\struck{soit $C$ un ensemble de} \add{dont le « noyau » est le système local} $\mathbb{Z}^{A(G)}$ sur $(\Pi M_{G},\Gamma')$, tordu grâce à l'action de $\Gamma'$ sur $A(G)$. Si $C \subset A(G)$ est stable par l'action de $\Gamma'$, alors on a un \uncertain{homomorphisme}
\[
(116)\qquad \mathbb{Z}^{A(G)} \longrightarrow \mathbb{Z}^{C}
\]

\page{85}\note{p. 42 de l'auteur.}
de systèmes locaux, qui permet de construire une extension quotient de $(S\Pi M_{G},\Gamma')$, extension du groupoïde $(\Pi M_{G},\Gamma')$ par $\mathbb{Z}^{C}$, qu'on va noter
\[
(117)\qquad (S_{C}\Pi M_{G},\Gamma') .
\]
\note{les lignes qui suivent, jusqu'à la formule (118), sont barrées de traits obliques.}
\struck{Dans le cas où $\Gamma' = \Gamma^{!}_{G,G'}$, avec $G' = G/C$, on \uncertain{retrouve} l'extension $\Pi M^{*}_{G,G'}$ de $\Pi M_{G,G'}$ par $\mathbb{Z}^{C}$ (tordu).}
\[
(118)\qquad \text{\struck{$(S_{C}M_{G},\Gamma_{G,G/C}) \simeq \Pi M^{*}_{G,G/C}$}}
\]
On va plutôt expliciter les choses en termes de $S_{A}\Pi M_{G} \to (\Pi M_{G},\Gamma^{!}_{G})$. Si $G'$ \add{$\simeq G/C$} \struck{est une} \struck{spécialisation} \add{généralisation} de $G$, alors on a
\[
(119)\qquad S_{A}\Pi M_{G} \longrightarrow S_{A'}\Pi M_{G'}
\]
qui passe aux quotients en
\note{diagramme (120) ; le numéro, surchargé, peut aussi se lire (119). Le terme en bas à gauche portait un indice biffé.}

\begin{tikzcd}
S_{A}\Pi M_{G} \arrow[r, "{\tau_{G,G'}}"] \arrow[d] & S_{A'}\Pi M_{G'} \arrow[d] \\
S_{C}\Pi M_{G} \arrow[r] & \Pi M_{G'}
\end{tikzcd}

compatible avec l'action de $\Gamma_{G,G'}$, d'où \struck{\ill{}} la deuxième flèche horizontale ; \uncertain{passant} \ill{}

\page{86}
et factorisation en
\[
(121)\qquad (S_{C}\Pi M_{G},\Gamma^{!}_{G,G'}) \longrightarrow \Pi M_{G'} ,
\]
\note{sous le premier membre de (121), une accolade porte $\Pi M^{*}_{G,G'}$, raturé.}
\struck{et} plus généralement, si $\Gamma^{!}_{G,G'} \subset \Gamma' \subset \Gamma_{G,G'}$ est un sous-groupe de $\Gamma = \mathrm{Aut}\, G$ qui stabilise $C$ et contient $\Gamma^{!}_{G,G'}$, on trouve
\[
(122)\qquad (S_{C}\Pi M_{G},\Gamma') \longrightarrow (\Pi M_{G'},\Gamma'/\Gamma^{!}_{G,G'})
\]
\[
\simeq (\Pi M_{G,G'},\Gamma'/\Gamma^{!}_{G,G'})
\]
\note{le signe $\simeq$ est écrit $\wr$ sous le second membre de (122).}

Ici la relation de transitivité (105) \uncertain{qu'on explique} s'exprime, plus simplement, par le fait que les \add{groupoïdes} $S_{A}\Pi M_{G}$ dépendent fonctoriellement des \add{MD-}graphes $G$, pour les morphismes de généralisation.

On va axiomatiser un peu \struck{plus}. \add{Les} MD-graphes \add{et morphismes de généralisation} forment une \emph{catégorie}, \supplied{où} tous les morphismes sont des épimorphismes. Cette catégorie satisfait
\marginal{$\mathcal{G}$ pour \uncertain{abréger}. La catégorie \uncertain{visée} d'un \uncertain{diagramme} de groupoïdes (\struck{\ill{}} $\mathcal{G}$) \uncertain{au-dessus}…}

\page{87}\note{p. 43 de l'auteur.}
à la condition suivante
\begin{quote}
(123) (Simpl) Pour tout $G \in \mathcal{G}$, l'ens.\ ordonné des quotients $G'$ de $G$ est \struck{\ill{}} \add{un « simplexe » i.e.\ \uncertain{isomorphe}} \struck{\ill{}} à l'ens.\ des parties d'un ens.\ fini.
\end{quote}
\note{dans (123), le symbole lu $\mathcal{G}$ ressemble à un $\mathcal{M}$.}
Si $A_{G}$ désigne l'ens.\ des quotients \add{minimaux} \struck{\ill{}} de $G$ qui sont $\neq G$ et qui \add{ordinaires} \supplied{sont} \ill{}, dans l'ens.\ des quotients de $G$ est \uncertain{donc} isomorphe à l'ens.\ des parties de $A_{G}$. (Si $C \subset A_{G}$, \ill{} désigne par $G/C$ le quotient de $G$ associé à $A_{G} \setminus C$ (\uncertain{d'autant} plus petit que $C$ est plus grand).)

On suppose donnés pour tout $G \in \mathrm{Ob}\,\mathcal{G}$ deux groupoïdes $S\Pi M_{G}$ et $\Pi M_{G}$, et un morphisme
\[
(124)\qquad S\Pi M_{G} \xrightarrow{\ \varphi_{G}\ } \Pi M_{G} .
\]
On suppose que $\varphi_{G}$ fait de $S\Pi M_{G}$ une extension de $\Pi M_{G}$ par $\mathbb{Z}^{A(G)}$ (en fait, c'est une structure supplémentaire sur $S\Pi M_{G}$). On suppose que $G \mapsto \Pi M_{G}$ est un foncteur en $G$ pour les \emph{iso.}\ (\uncertain{resp.} pour les morphismes quelc.)
\marginal{NB Il faudrait \uncertain{montrer} comment cette axiomatisation s'applique \add{à la situation des multiplicités} \struck{topologiques relatives \ill{}} \uncertain{définies} par un diviseur à croisements normaux… (cela \ill{} \ill{}). \uncertain{Noter} que $S\Pi M_{G}$ est \uncertain{une} extension de $\Pi M_{G}$ \uncertain{non pas} par un $\mathbb{Z}^{C}$ tordu, mais \uncertain{quotient} d'une \uncertain{telle}…}

\page{88}
et que \add{$G \mapsto$} $S\Pi M_{G}$ est un foncteur en \struck{\ill{}} $G$ pour les morphismes quelconques, et que (124) est fonctoriel pour les iso. Pour un morphisme quelconque $G \xrightarrow{f} G'$, correspondant à une partie $C \subset A_{G}$ (telle que $G \to G'$ soit isom.\ à $G \to G/C$), \struck{\ill{}} quotient de $G$, on suppose que l'on a un diagramme commutatif

\begin{tikzcd}
S\Pi M_{G} \arrow[r, "{S\Pi(f)}"] \arrow[d, "{\varphi_{G}}"'] & S\Pi M_{G'} \arrow[d, "{\varphi_{G'}}"] \\
S_{C}M_{G} \arrow[r, "{\gamma_{f}}"] & \Pi M_{G'}
\end{tikzcd}

\note{diagramme (125) ; le $\Pi$ manque dans $S_{C}M_{G}$, ici comme deux lignes plus bas, et l'indice $C$ est écrit sur un $A$ biffé.}
i.e.\ que le composé
\[
S\Pi M_{G} \xrightarrow{\ S\Pi(f)\ } S\Pi M_{G'} \xrightarrow{\ \varphi_{G'}\ } \Pi M_{G'}
\]
se factorise par l'extension $S_{C}M_{G}$ de $\Pi M_{G}$ par $\mathbb{Z}^{C}$, déduite de l'extension $S\Pi M_{G}$ de $\Pi M_{G}$ par $\mathbb{Z}^{A(G)}$ grâce à l'homom.\ canonique $\mathbb{Z}^{A(G)} \to \mathbb{Z}^{C}$. On a automatiquement la relation de transitivité pour un composé.

\page{89}\note{p. 44 de l'auteur.}
\section{10) Structure MDT discrète : étude des morphismes de normalisation entre MD-graphes}
\note{titre souligné par l'auteur ; « MD- » est ajouté au-dessus de « graphes ».}

Nous voulons expliciter, en plus des morphismes \struck{$\varphi_{G}$} $S\Pi f$ ($f$ morphisme de généralisation) et $\Pi f$ ($f$ iso de MD-graphes) \uncertain{entre autres} des morphismes du type
\[
M_{G} \longrightarrow M_{g_{\alpha},\hat{I}_{\alpha}}
\]
\uncertain{etc.}
\[
M^{\natural}_{g,I} \longrightarrow M_{g,I'} \quad \text{pour } I' \subset I .
\]
\note{l'exposant de $M_{g,I}$, rendu $\natural$, est douteux ; l'indice $I$ est surchargé.}

On va introduire pour ceci une notion de morphismes de MD-graphes, qu'on appellera morphismes de normalisation (l'intersection des deux notions de morphismes est celle d'iso de MD-graphes). Géométriquement, elle correspond à la situation d'une courbe stable \struck{$(\hat{X},I)$} \add{$X = \hat{X} \setminus I$} sur un corps alg.\ clos $k$, où on distingue une partie $I' \subset I$ de l'ens.\ des pts à l'infini, et une partie $A' \subset A$ de l'ens.\ des pts doubles de $X$,

\page{90}
et en considérant la courbe $X'$, ayant $I'$ comme ensemble de pts à l'infini, et $A'$ comme ens.\ des pts doubles, obtenue en « bouchant les trous » relatifs aux \uncertain{pts} de $I \setminus I'$, \struck{\ill{}} et en normalisant $X$ en les pts de $A \setminus A'$. On trouve donc $X' = \hat{X}' \setminus I'$, $\hat{X}'$ étant le normalisé de $\hat{X}$ relatif aux pts de $A \setminus A'$. On fera attention que $\hat{X}'$ ne \emph{s'envoie pas} dans $\hat{X}$ (sauf si $I' = I$), mais seulement dans $\hat{X}$.
\note{les accents des deux premiers $\hat{X}$ de cette phrase sont retouchés ; le sens appelle plutôt « $X'$ ne s'envoie pas dans $X$ ».}
\struck{On peut aussi \ill{} plus loin, et} La construction de $X'$ en termes de $X$ est donc complète, dès qu'on se donne des parties $I' \subset I$ et $A' \subset A$ -- cela revient, en termes du \struck{\ill{}} graphe associé à $X$, à se donner
\begin{quote}
a) un sous-graphe $G'$ de $G$ ayant même ensemble de sommets ;

b) un sous-ens.\ $I'$ de $I$.
\end{quote}
\struck{Mais on} Le graphe de $X'$ est alors le sous-graphe en question. Le \struck{MD}\uncertain{-graphe}

\page{91}\note{p. 45 de l'auteur.}
complète de $X'$ se déduit de celle de $X$, en prenant la \uncertain{même} \uncertain{application genre}, et $I' \to S$ induite par $I \to S$.

En fait, quand on a normalisé $\hat{X}$ en les pts de $A \setminus A'$ pour obtenir $\hat{X}'$, on a sur $\hat{X}'$, en plus des pts marqués indexés par $I$ et des pts \add{doubles} \struck{\ill{}} indexés par $A'$, un \uncertain{deuxième} ens.\ de « pts marqués », savoir l'image inverse \add{$\tilde{J}$} de $J = A \setminus A'$, \struck{E} de cardinal égal à $2\,\mathrm{Card}(A \setminus A')$. En termes de la courbe normalisée propre $Y$ \struck{\ill{}} contenant $\tilde{A} \amalg I$, qui définit $\hat{X}$ comme $Y/\sigma_{\tilde{A}}$ et $X$ comme $(Y \setminus I)/\sigma_{\tilde{A}}$, $\hat{X}'$ est défini comme $Y/\sigma_{\tilde{A}'}$, ce qui nous donne un ens.\ de pts marqués \struck{\ill{}} $\tilde{J} \amalg I$ sur $\hat{X}'$. Ceci étant, au lieu d'« oublier » $\tilde{J}$ et $I \setminus I'$ pour ne retenir que l'ens.\ de pts marqués

\page{92}
$I'$ sur $\hat{X}'$, on peut aussi, plus généralement, se donner une partie $I'$ (non pas de $I$, mais de) $\tilde{J} \amalg I$, qu'on retiendra comme « pts marqués ». Quand on s'est ainsi donné (pour un \struck{\ill{}} MD-graphe $\underline{G}$ donné) une partie $A'$ de $A(G)$, et une partie $I'$ de $I \amalg \tilde{J}$ (où $\tilde{J} = \tilde{A}(G) - \tilde{A}'$), alors pour toute courbe spéciale $G$-épinglée, on déduit une courbe spéciale, qui sera épinglée par le MD-graphe $\underline{G}'$ déduit de $\underline{G}$, $A'$, $I'$ en prenant le sous-graphe \add{$G'$ de $G$} \uncertain{ayant} mêmes sommets et ayant $A'$ comme ens.\ d'arêtes, \uncertain{en prenant} comme « ens.\ de pts marqués » $I' \subset I \amalg \tilde{J}$, qui s'envoie dans $S$ \struck{\ill{}} par
\[
(126)\qquad I' \longrightarrow S \quad \text{induite par } I \to S \text{ et } \tilde{A}(G) \xrightarrow{\ \text{origine}\ } S ,
\]
l'application genre pour $\underline{G}'$ étant celle pour $\underline{G}$.

\page{93}\note{p. 46 de l'auteur.}
On peut généraliser encore un peu cette construction, en \struck{choisissant \ill{}} \add{remplaçant $X'$ par} une réunion de composantes connexes \struck{de composantes irréductibles}, contenant $A'$ et $I'$. Cela revient à prendre un sous-graphe \emph{quelconque} $G'$ de $G$ (n'ayant pas nécessairement \struck{\ill{}} même ens.\ de sommets, mais un ens.\ de sommets $S' \subset S$), [défini donc par la donnée d'une partie $\tilde{A}'$ de $\tilde{A}(G)$ stable par $\sigma_{\tilde{A}}$ et telle que $o_{G} : \tilde{A} \to S$ applique $\tilde{A}'$ dans $S'$], plus une partie $I'$ de $I \amalg \tilde{J}$ (où $\tilde{J} = \tilde{A} \setminus \tilde{A}'$), appliquée dans $S'$ par l'\struck{application origine}
\[
I \amalg \tilde{J} \longrightarrow S \quad \text{induite par } I \to S \text{ et par } \tilde{A} \xrightarrow{\ o_{G}\ } S .
\]

Un cas particulièrement important est celui où $G'$ est un sous-graphe ponctuel, \add{donc $\tilde{J} = \tilde{A}$,} défini par un $\alpha \in S$, et où on prend
\[
I' = I_{\alpha} \amalg \tilde{A}_{\alpha} = \text{image inverse de } \{\alpha\} \text{ par } I \amalg \tilde{A} \to S .
\]

Si on veut exprimer ces constructions en termes d'une notion adéquate de

\page{94}
« morphisme » (de normalisation) de MD-graphes, on aboutit à la

\emph{Définition}. Soient
\[
\underline{G} = (\underbrace{S,\tilde{A},\sigma_{\tilde{A}},o_{G}}_{G},\ I \to S \xrightarrow{\ g\ } \mathbb{N})
\quad\text{et}\quad
\underline{G}' = (\underbrace{S',\dots}_{G'})
\]
deux MD-graphes. Un \emph{morphisme de \struck{\ill{}} normalisation} de $\underline{G}'$ dans $\underline{G}$ est la donnée d'un couple $(f_{G},f_{I})$, où
\note{le mot qui précède « morphisme » est écrit sur une rature, lu « Un » avec doute.}
\begin{quote}
a) $f_{G} : G' \to G$ est un homomorphisme de graphes, qui est injectif sur les sommets et injectif sur l'ens.\ des arêtes (donc identifie $G'$ à un sous-graphe de $G$) ;

b) $f_{I} : I' \to I \amalg \tilde{A}$ \struck{\ill{}} est une application injective, dont l'image est contenue dans $I \amalg \tilde{J}$, où $\tilde{J} = \tilde{A}(G) - f_{G}(\tilde{A}')$, et rendant commutatif le diagramme
\end{quote}

\begin{tikzcd}
I' \arrow[r, hook] \arrow[d] & I \amalg \tilde{A} \arrow[d] \\
S' \arrow[r, hook] & S
\end{tikzcd}

\note{diagramme (129) ; la flèche du bas porte une étiquette biffée.}
\marginal{On a supposé de plus c) $\forall s' \in S'$, $g'(s') = g(f_{G}(s'))$.}
\marginal{On dit que $f$ est un \emph{morphisme de normalisation strict} si $S' \xrightarrow{\sim} S$ et si $I' \xrightarrow{\sim} I \amalg \tilde{J}$.}

Par \uncertain{composition}, \uncertain{catégorie} des MD-graphes et morphismes de normalisation. \struck{\ill{}}

Voici une façon de raccorder cette notion à une notion « orthodoxe » de morphismes de graphes, en mettant « dans un même sac », pour un MD-graphe

\page{95}\note{p. 47 de l'auteur.}
donné $G$, l'ens.\ $\tilde{A}(G)$ et $I(G)$ (munis l'un et l'autre d'une application \add{« origine »} à valeurs dans $S(G)$). Pour ceci, on va utiliser la notion de « graphe avec boucles aplaties », qui généralise celle de graphe, et topologiquement correspond à la classification des couples $(X,S)$ d'un espace topologique compact $X$, et d'une partie finie $S$, tels que a) $X$ soit \uncertain{étoilé} au voisinage de tout pt de $S$, et b) toute composante connexe de $X \setminus S$ est une 1-variété \uncertain{à bord} non compacte (on \uncertain{veut} que toute \struck{\uncertain{adhérence}} d'« arête » \uncertain{aboutisse} en un sommet). Il revient au même de dire que $X$ est un 1-complexe top., et $S$ une partie finie qui contient les sommets d'ordre $\geq 3$, et telle qu'aucune composante connexe de $X \setminus S$ ne soit composante connexe de $X$ (i.e.\ que toute composante connexe de $X$ contienne un sommet, i.e.\ $S \to \pi_{0}(X)$ est surjectif).

Ceci se classifie combinatoirement par des systèmes $(S,\hat{\tilde{A}},\sigma_{\hat{\tilde{A}}},o_{\tilde{A}})$, où \struck{$\tilde{A}$, \ill{}} $S$, $\hat{\tilde{A}}$ sont des ens.\ finis, $\sigma_{\tilde{A}}$ une involution de $\tilde{A}$, et $o_{\tilde{A}} : \tilde{A} \to S$ une
\note{l'auteur hésite entre $\tilde{A}$ et $\hat{\tilde{A}}$ (accents superposés) dans ce paragraphe ; on transcrit tel quel.}

\page{96}
application. On n'exige plus que $\sigma_{\hat{A}}$ soit sans pts fixes -- l'ens.\ des pts fixes, $I$, \struck{se} correspond aux « boucles aplaties » du \add{1-}complexe, i.e.\ aux \struck{« \ill{} »} « pts-bouts » de $X$ qui ne sont pas dans $S$. La donnée d'un graphe à boucles aplaties revient donc (en désignant par $\hat{I}$ \uncertain{ou} $I \amalg \tilde{A}$, où $I = \hat{I}^{\sigma}$ et $\tilde{A} = \hat{I} \setminus \hat{I}^{\sigma}$) à la donnée $(S,\tilde{A},\sigma_{\tilde{A}},o_{\tilde{A}})$ d'un \struck{\ill{}} graphe ordinaire, plus d'un ensemble $I$ et d'une application $I \to S$. On décrit un tel graphe \add{à boucles aplaties} par un graphisme tel que

\drawing{un graphe : un sommet de degré trois porte deux boucles aplaties (segments libres) et est relié à un second sommet par deux arêtes parallèles ; ce second sommet est relié par une arête à un troisième sommet qui porte une boucle ; un quatrième sommet, en bas, relié au premier, porte trois boucles aplaties.}

où les sommets sont figurés par des points gros $\bullet$, et les boucles aplaties \struck{\ill{}} se distinguent par le fait justement qu'elles n'ont pas de point gros. Le graphe \add{à b.a.} ci-dessus correspond donc, en notation des « pts marqués », à

\drawing{le même dessin réduit à sa partie supérieure : un sommet portant deux boucles aplaties numérotées 2 et 3, avec le numéro 1 au-dessus du sommet ; deux arêtes parallèles vers un second sommet, puis une arête vers un sommet portant une boucle. La phrase s'interrompt après le dessin.}

\page{97}\note{p. 48 de l'auteur.}
\marginal{cas d'une \uncertain{normalisation} stricte.}
Dans cette optique, on voit que dans le cas où \add{$S' \simeq S$,} $I' \simeq I \amalg \tilde{J}$, on obtient $G'$ à partir de $G$ en se donnant l'ens.\ $C$ des \struck{flèches} \add{arêtes (ordinaires)} de $G$ « qu'on veut couper » (complémentaire de $I'$ \uncertain{dans} $\hat{I}' \to \hat{I}$) en prenant un pt sur chaque arête $\in C$, et en « coupant » l'arête en ce point avec un coup de ciseaux (i.e.\ prenant le \uncertain{diagramme} de $X$ suivant \uncertain{la partie finie marquée en} $X$)

\drawing{un graphe : un sommet trivalent portant deux boucles aplaties en haut, relié à droite par une arête marquée d'une croix $\times$ à un second sommet, lui-même relié à un troisième sommet portant une boucle, laquelle est marquée d'une croix ; un quatrième sommet, en bas à gauche, relié au premier.}

donne

\drawing{le graphe $G'$ obtenu en coupant aux croix : le sommet trivalent, l'arête coupée devenue deux boucles aplaties, la boucle coupée devenue deux boucles aplaties issues du troisième sommet.}

On trouve un nouveau graphe à boucles aplaties (qui a \uncertain{toutes les chances} d'être disconnecté, \uncertain{même} si $G$ était connexe ; si $G$ \uncertain{n'avait} pas de boucles aplaties, il a \uncertain{tendance à en avoir}…). Si de plus on remplace $G'$ par un sous-graphe de $G'$ en enlevant certaines boucles aplaties, et en prenant une réunion de composantes connexes de ce qui reste (\struck{\ill{}} \add{\ill{}}) c'est
\marginal{En somme, le \uncertain{graphe} déduit de $G$ en donnant une \uncertain{involution} \ill{} pts fixes de l'ens.\ de ses bouts d'« arêtes », ou « recollés » \uncertain{celles-ci} \ill{} \uncertain{pour} \ill{} bouts.}

\page{98}
kif-kif, à la composante connexe près (en ajoutant ou retranchant des boucles aplaties), on trouve la relation la plus générale entre $G$ et un \uncertain{$G'$}, des \uncertain{types} explicités dans la définition des « morphismes de normalisation ».

On définit de façon évidente les composés de morphismes de normalisation. Le composé de morphismes de normalisation strict est strict. Les MD-graphes \uncertain{et} morphismes de normalisation stricts forment une catégorie dont les iso.\ sont les iso.\ ordinaires de MD-graphes. Les \uncertain{morphismes} de cette catégorie sont \uncertain{des} monomorphismes. \struck{\ill{}} \add{\uncertain{Quand}} on se borne aux morphismes \add{\uncertain{dits}} stricts, on voit que (pour les morphismes de normalisation \uncertain{généraux}) l'ens.\ des sous-MD-graphes $G'/\!\simeq$ \uncertain{stricts} d'un MD-graphe $G$ (\struck{\ill{}} \add{\uncertain{classes}} \ill{}) \uncertain{s'identifie} à l'ens.\ des parties de $A(G)$, en associant à une partie $C \subset A(G)$

\page{99}\note{p. 49 de l'auteur.}
le « \emph{diagramme} » de $G$ suivant $C$, obtenu en « coupant » comme il a été dit les \uncertain{arêtes} ou les milieux des arêtes $\in C$ (ce qui revient à « enlever » les arêtes, et à remplacer les \uncertain{sommets} correspondants par des pts marqués.)

Soit un morphisme de normalisation entre MD-graphes
\[
(130)\qquad G' \longrightarrow G .
\]
Si c'est un morphisme de normalisation \emph{strict}, on voit que $G'$ est stable \uncertain{ssi} $G$ l'est. (Dans le cas général, on voit que si $S' \to S$ est surjectif, et si $G'$ est stable, $G$ l'est aussi.)

\struck{Supposons} Si $X$ est une courbe spéciale (sur une base quelconque $S$), épinglée par $G$, \uncertain{donc} on voit (en procédant comme au début de la section) qu'on lui associe canoniquement une courbe spéciale $C'$ épinglée, donc un foncteur
\[
(131)\qquad M_{G}(S) \longrightarrow M_{G'}(S) ,
\]
\note{le numéro (131) est écrit sur un autre chiffre.}
qui pour $S$ variable est compatible avec changement de base, et donc définit un morphisme de multiplicités modulaires (si $G$, $G'$ sont stables)

\page{100}
\[
(132)\qquad M_{G} \longrightarrow M_{G'} ,
\]
\uncertain{avec} \uncertain{transitivité} évidente pour des composés. Ainsi $M_{G}$ est un \emph{foncteur contravariant en $G$} pour les \add{morphismes de normalisation \uncertain{stables}} \struck{morphismes de normalisation}. Ceci posé, on a

\emph{Proposition}. a) Si $G \to G'$ est un morphisme de normalisation \emph{strict}, alors (132) est un \emph{iso.}\ de multiplicités.

b) Soient $G_{\alpha}$ des \add{MD-graphes stables} \struck{graphes connexes}, et $G$ le MD-graphe somme des $G_{\alpha}$. Alors le morphisme
\[
(133)\qquad \text{\struck{$\Pi_{1}M_{G}$, $\Pi_{1}$}}\ M_{G} \longrightarrow \prod_{\alpha} M_{G_{\alpha}}
\]
est un iso.\ de multiplicités.
\marginal{NB On \uncertain{remplace} \uncertain{(133)} \uncertain{par} \uncertain{une} \uncertain{variante} $\Pi M_{G}$ \uncertain{pour} \uncertain{un} \ill{} \uncertain{en produit}.}

\emph{Cor.} $\Pi_{1}M_{G}$ est un foncteur contravariant en le MD-graphe stable $G$, \struck{pour} les
\note{la phrase se poursuit au-delà de ce lot.}

\end{document}
